\documentclass[a4paper,11pt]{article}

\usepackage[margin=2.5cm]{geometry}
\usepackage[T1]{fontenc}
\usepackage[utf8]{inputenc}
\usepackage{listings}
\usepackage{hyperref}

\title{Practical Work 4: Word Count}
\author{Nguyen Trong Bach \\ Student ID: 23BI14057}
\date{\today}

\begin{document}

\maketitle

\section{Introduction}
This practical implements the classical Word Count example using
a simple MapReduce style framework written in C++. Instead of using
a full MapReduce system such as Hadoop, I design a minimal
implementation with clearly separated \textbf{Mapper} and
\textbf{Reducer} functions. The goal is to count how many times each
word appears in an input text file.

\section{Choice of Implementation}
C++ is chosen for this practical because:
\begin{itemize}
  \item it is close to C and fits well with the previous practical works;
  \item the Standard Library provides convenient containers such as
  \texttt{std::map} to store key--value pairs;
  \item the program can be compiled and executed easily on the course
  environment using \texttt{g++}.
\end{itemize}

\section{Mapper and Reducer Design}

\subsection*{Mapper}
The Mapper is implemented by the function
\texttt{map\_line(const std::string\&, std::map<std::string,int>\&)}.
For each input line the Mapper:
\begin{enumerate}
  \item scans the line character by character;
  \item groups sequences of alphabetic characters into words;
  \item converts each word to lowercase to avoid case differences;
  \item increases the local counter for that word in a local
  map of type \texttt{std::map<std::string,int>}.
\end{enumerate}
In logical MapReduce terms, for every word the pair
\texttt{(word, 1)} is produced and stored in the local map.

\subsection*{Reducer}
The Reducer is implemented by the function
\texttt{reduce\_counts(const std::map<std::string,int>\&,
std::map<std::string,int>\&)}. The Reducer receives a local map
containing pairs of the form \texttt{(word, count)} and merges them
into a global map:
\begin{itemize}
  \item for each key in the local map the corresponding value is added
  to the value stored in the global map;
  \item if the key does not exist in the global map yet, it is inserted
  with the given count.
\end{itemize}
This behaviour is similar to the Reduce phase in MapReduce, where all
values associated with the same key are combined.

\section{System Organization}
The whole Word Count system is contained in a single source file
\texttt{word\_count.cpp}. The main program follows these steps:
\begin{enumerate}
  \item read the input file name and output file name from the command
  line arguments;
  \item open the input file and read it line by line;
  \item for each line, call the Mapper to obtain local word counts;
  \item call the Reducer to merge the local counts into a global map;
  \item after all lines have been processed, write the global word
  counts to the output file, one word per line.
\end{enumerate}
The output format is very simple: each line contains a word followed by
its total count, for example:
\begin{verbatim}
apple  5
banana 3
\end{verbatim}

\section{Implementation Snippet}
The following listing shows the core part of the Mapper function:

\begin{lstlisting}[language=C++]
void map_line(const std::string &line,
              std::map<std::string, int> &local_counts) {
    std::string word;
    for (char c : line) {
        if (std::isalpha(static_cast<unsigned char>(c))) {
            word.push_back(
                static_cast<char>(std::tolower(
                    static_cast<unsigned char>(c))));
        } else {
            if (!word.empty()) {
                ++local_counts[word];
                word.clear();
            }
        }
    }
    if (!word.empty()) {
        ++local_counts[word];
    }
}
\end{lstlisting}

The Reducer simply iterates over the local map and adds all counts to
the global map.

\section{Group Work}
This practical work was completed individually. I was responsible for
designing the simple MapReduce style, implementing the C++ program and
writing this report.

\section{Conclusion}
This practical demonstrates that the MapReduce programming model can be
implemented even without a full framework. By separating the logic into
Mapper and Reducer functions, the program mirrors the core ideas of
MapReduce while remaining small and easy to understand. The same design
could later be extended to run in parallel using multiple processes or
machines.

\end{document}
